\begin{tabular}{llrlrlr}
\toprule
\textbf{Punto} & \textbf{Fuente} & \textbf{$D$ (sesion valida)} & \textbf{Medida} & \textbf{$D$ (tierra mal)} & \textbf{Medida} & \textbf{Cambio [\%]} \\
\midrule
D @ 1600 V & Cs-137 & $0.4934 \pm 0.0011$ & 2026-09-08 & $0.3599 \pm 0.0004$ & 2026-09-01 & $-27.1$ \\
D @ 1650 V & Cs-137 & $0.6323 \pm 0.0010$ & 2026-09-08 & $0.4668 \pm 0.0010$ & 2026-09-01 & $-26.2$ \\
D @ 1700 V & Cs-137 & $0.7877 \pm 0.0019$ & 2026-09-08 & $0.6171 \pm 0.0010$ & 2026-09-01 & $-21.7$ \\
\bottomrule
\end{tabular}
\par\smallskip
\begin{minipage}{0.95\linewidth}\footnotesize Mismo detector, mismo voltaje nominal y misma fuente. La sesion del 1-sep se tomo con la tierra del detector mal conectada y esta descartada; la del 8-sep es la valida. La columna del cambio mide, entonces, cuanto movio la ganancia esa falla, no la reproducibilidad del montaje.\end{minipage}
